\documentclass[11pt, english, letterpaper]{article}
\usepackage[T1]{fontenc}
\usepackage[utf8]{inputenc}
\usepackage{textcomp}
\usepackage{lmodern}
\usepackage{graphicx}
\usepackage[usenames,dvipsnames,table]{xcolor}
% xcolor must come before pdfpages, which loads it with no options.
\usepackage{pdfpages}
\usepackage[margin=0.9in]{geometry}
\usepackage{babel}
\usepackage{caption}
\usepackage{float}
\usepackage{longtable}
\usepackage{booktabs}
\usepackage{tabularx}
\usepackage[shortlabels]{enumitem}
\usepackage{url}
\usepackage[colorlinks=true, linkcolor=black, urlcolor=black, citecolor=black,
linkbordercolor={white}, anchorcolor=black]{hyperref}
%
% Glyphs the Python exhibits emit that pdflatex does not map by default.
% \ensuremath is used instead of dollar-sign math: Stata would try to expand
% anything after a dollar sign as a global macro.
\DeclareUnicodeCharacter{2265}{\ensuremath{\geq}}
\DeclareUnicodeCharacter{2264}{\ensuremath{\leq}}
\DeclareUnicodeCharacter{00D7}{\ensuremath{\times}}
\DeclareUnicodeCharacter{2212}{--}
\DeclareUnicodeCharacter{2013}{--}
\DeclareUnicodeCharacter{2014}{---}
\DeclareUnicodeCharacter{00B0}{\textdegree}
\DeclareUnicodeCharacter{00B2}{\textsuperscript{2}}
\DeclareUnicodeCharacter{2018}{\textquoteleft}
\DeclareUnicodeCharacter{2019}{\textquoteright}
\DeclareUnicodeCharacter{201C}{\textquotedblleft}
\DeclareUnicodeCharacter{201D}{\textquotedblright}
\DeclareUnicodeCharacter{03C3}{\ensuremath{\sigma}}
%
\pdfminorversion=6
\linespread{1.15}
\def\arraystretch{1.25}
\captionsetup{font=small, labelfont=bf, justification=raggedright, singlelinecheck=false}
\setlength{\parskip}{6pt}
\setlength{\parindent}{0pt}
%
% Exhibits live beside this .tex, so the folder compiles anywhere.
\newcommand*{\figTwo}{exhibits}%
%
\begin{document}
\title{Winter Weather Panel: Descriptive Report\\[4pt]
\large PRISM county--month data, 1981--2025 \\[2pt]
\normalsize Version 2}
\date{-2026-09-11}
\maketitle
\thispagestyle{empty}
\vspace{-1em}
\begin{center}\begin{minipage}{0.86\textwidth}\small
\textbf{Provenance.} All exhibits generated by
\path{codePYTHON/09_descriptive_weather_full.ipynb}, executed on Kodama against the
complete PRISM county-month panel. Figures are vector PDFs reproduced without
rescaling of content. Tier~1 QA exhibits are indexed in Appendix~B rather than
reproduced.
\par\vspace{6pt}
\textbf{What changed since the 2026-08-28 report.} Titles and source notes have been
removed from the images and the notes now appear beneath each figure; map legends are
horizontal and sit at lower left; state borders are drawn in black. New exhibits: all
state trends pooled into one figure, a Data-section composite, a full-sample mean
winter temperature map, an early/recent/change treatment of extreme-cold days, the
Winter Severity Index, area-weighted companions to the national exhibits, and the
cross-season anomaly and trend comparisons.
\par\vspace{6pt}
\textbf{One definition changed.} The warm-anomaly exhibits previously flagged a
\emph{detrended} residual above 1.5 SD. They now use the same construction as
\path{warm_winter_1sd} and \path{warm_winter_2sd} in the merged county-year panel:
the raw winter mean above the county's own full-sample mean by one or two standard
deviations. The figures and the regression table therefore describe the same object.
Section~\ref{sec:decisions} lists this and the remaining choices requiring a decision.
\end{minipage}\end{center}
\newpage
\tableofcontents
\newpage
\section{Coverage and data integrity}
The panel is complete: no duplicate keys, no unmatched rows across the monthly
and derived files, no incomplete county-months, and no QA-field mismatches. The
excluded county-winters are exactly the two boundary winters that cannot be
complete, not unexplained missingness.
\par
\begin{longtable}{@{}lr@{}}
\caption{Coverage and integrity diagnostics}\\
\toprule Metric & Value \\ \midrule \endfirsthead
\toprule Metric & Value \\ \midrule \endhead
\bottomrule \endfoot
dataset & PRISM \\
monthly rows & 1678320 \\
derived rows & 1678320 \\
merged rows & 1678320 \\
counties & 3108 \\
first year & 1981 \\
last year & 2025 \\
duplicate monthly keys & 0 \\
duplicate derived keys & 0 \\
monthly only rows & 0 \\
derived only rows & 0 \\
incomplete months & 0 \\
qa mismatch n days & 0 \\
qa mismatch ppt total & 0 \\
qa mismatch dataset types & 0 \\
qa mismatch expected days & 0 \\
qa mismatch is incomplete & 0 \\
candidate county winters & 142968 \\
complete county winters & 136752 \\
excluded incomplete county winters & 6216 \\
first complete winter & 1982 \\
last complete winter & 2025 \\
expected calendar years & 1981–2025 \\
observed calendar years & 1981–2025 \\
observed complete winters & 1982–2025 \\
\end{longtable}
\begin{longtable}{@{}lrl@{}}
\caption{Tier 1 QA findings}\\
\toprule Check & Result & Status \\ \midrule \endfirsthead
\toprule Check & Result & Status \\ \midrule \endhead
\bottomrule \endfoot
Duplicate county-year-month keys & 0 & PASS \\
Rows appearing in only one merge source & 0 & PASS \\
Incomplete county-months & 0 & PASS \\
Excluded incomplete county-winters & 6216 & INFO \\
County-variable mean flags (|z| ≥ 5) & 20 & REVIEW \\
\end{longtable}
\clearpage
\section{Long-run winter warming}
\begin{figure}[H]
\centering
\includegraphics[width=\textwidth]{\figTwo/prism_national_long_run_warming.pdf}
\caption{National mean winter temperature, 1982--2025, with a fitted linear trend.}
\end{figure}
\vspace{-6pt}
{\footnotesize\textit{Notes.} Source: PRISM; complete DJF winters only. Estimand: the average county (counties weighted equally), not an area-, population- or road-weighted national average.\par}
\clearpage
\begin{figure}[H]
\centering
\includegraphics[width=\textwidth]{\figTwo/prism_national_long_run_warming_area_weighted.pdf}
\caption{The same series with counties weighted by land area. The two answer different questions and both are reported.}
\end{figure}
\vspace{-6pt}
{\footnotesize\textit{Notes.} Source: PRISM; complete DJF winters only. Estimand: the average square kilometre of the contiguous United States (counties weighted by land area). Eyal ruled out population weighting (8/28); counties in the west are much larger than counties in the east, so the unweighted and area-weighted series answer different questions.\par}
\clearpage
\begin{figure}[H]
\centering
\includegraphics[width=\textwidth]{\figTwo/prism_state_trends_pooled.pdf}
\caption{Every state's fitted winter-temperature trend in one figure. States differ in both level and slope.}
\end{figure}
\vspace{-6pt}
{\footnotesize\textit{Notes.} Source: PRISM; complete DJF winters. One fitted linear trend per state (48 contiguous states plus DC), drawn in a single colour with transparency — the figure shows the spread of intercepts and slopes, and is not meant to make any individual state legible. Labelled: the three steepest warming trends (VT (+0.54), ME (+0.50), NH (+0.48)) and the three shallowest (ND (+0.05), SD (+0.06), MT (+0.09)), in °C per decade.\par}
\clearpage
\begin{figure}[H]
\centering
\includegraphics[width=\textwidth]{\figTwo/prism_state_trends_pooled_by_region.pdf}
\caption{The same trends coloured by EPA climate-impact region.}
\end{figure}
\vspace{-6pt}
{\footnotesize\textit{Notes.} Source: PRISM; the same fitted state trends coloured by EPA climate-impact region. Regions are assigned whole states, so no state is split; EPA Level I ecoregions cut across state lines and would need a county-level assignment.\par}
\clearpage
State-level trends are reproduced panel by panel in Appendix~\ref{app:states}.
\clearpage
\section{Where winters have warmed}
\begin{figure}[H]
\centering
\includegraphics[width=\textwidth]{\figTwo/prism_winter_warming_data_section.pdf}
\caption{Proposed Data-section figure: the change between the two periods above, the two period means below.}
\end{figure}
\vspace{-6pt}
{\footnotesize\textit{Notes.} Source: PRISM; complete DJF winters. Top: the change in mean winter temperature between the two periods, on a diverging scale centred at zero. Bottom: the two period means on a shared scale. Each county value is a mean over the 10 winters in its period (1982–1991 and 2016–2025), not a comparison of single winters or of sample endpoints. County borders are suppressed; state borders are drawn in black.\par}
\clearpage
\begin{figure}[H]
\centering
\includegraphics[width=\textwidth]{\figTwo/prism_mean_winter_temperature.pdf}
\caption{Mean winter temperature over the full sample -- the climatology the anomaly measures are taken against.}
\end{figure}
\vspace{-6pt}
{\footnotesize\textit{Notes.} Source: PRISM; mean of each county's complete DJF winters, 1982–2025. This is the climatology each county's warm-winter anomalies are measured against, so it is the baseline the anomaly maps standardise away.\par}
\clearpage
\begin{figure}[H]
\centering
\includegraphics[width=\textwidth]{\figTwo/prism_winter_temperature_early_vs_recent.pdf}
\caption{Mean winter temperature in the early and recent decades, on a shared colour scale.}
\end{figure}
\vspace{-6pt}
{\footnotesize\textit{Notes.} Source: PRISM; complete DJF winters. Both panels share one colour scale, so they are directly comparable. County borders are suppressed; state borders are drawn in black. Each county value is a mean over the 10 winters in its period (1982–1991 and 2016–2025), not a comparison of single winters or of sample endpoints.\par}
\clearpage
\begin{figure}[H]
\centering
\includegraphics[width=\textwidth]{\figTwo/prism_winter_temperature_change.pdf}
\caption{Change in mean winter temperature between the two periods.}
\end{figure}
\vspace{-6pt}
{\footnotesize\textit{Notes.} Source: PRISM; complete DJF winters. Red is warming, blue is cooling, on a diverging scale centred at zero. The comparison is 2016–2025 minus 1982–1991. Each county value is a mean over the 10 winters in its period (1982–1991 and 2016–2025), not a comparison of single winters or of sample endpoints.\par}
\clearpage
\section{Winter cold extremes}
\begin{figure}[H]
\centering
\includegraphics[width=\textwidth]{\figTwo/prism_extreme_cold_day_count.pdf}
\caption{Mean count of winter days with a daily minimum below 0\textdegree F.}
\end{figure}
\vspace{-6pt}
{\footnotesize\textit{Notes.} Source: PRISM; days with a daily minimum below 0°F, averaged over complete DJF winters. The 0°F threshold follows the metabolic-stress cutoff in the wildlife literature, not the 32°F road-conditions threshold used elsewhere in the panel. This is NOT the Winter Severity Index: the WSI adds a snow-depth hazard component and runs December through April. See the Winter Severity Index exhibit for that.\par}
\clearpage
\begin{figure}[H]
\centering
\includegraphics[width=\textwidth]{\figTwo/prism_extreme_cold_early_recent_change.pdf}
\caption{Extreme-cold days in the early and recent decades, and the change between them.}
\end{figure}
\vspace{-6pt}
{\footnotesize\textit{Notes.} Source: PRISM; days with a daily minimum below 0°F, averaged over the complete DJF winters in each period. Top panels share one scale. In the bottom panel RED marks counties that now get FEWER extreme-cold days — the colour scale is deliberately reversed relative to the temperature-change map, because for a count of cold days a decrease is the warming signal. Each county value is a mean over the 10 winters in its period (1982–1991 and 2016–2025), not a comparison of single winters or of sample endpoints.\par}
\clearpage
\section{Winter Severity Index}
\begin{figure}[H]
\centering
\includegraphics[width=\textwidth]{\figTwo/prism_winter_severity_index.pdf}
\caption{Mean Winter Severity Index (Kohn 1975 / Wisconsin DNR), read from the merged county-year panel.}
\end{figure}
\vspace{-6pt}
{\footnotesize\textit{Notes.} Winter Severity Index (Kohn 1975 / Wisconsin DNR): the count of days from 1 December to 30 April with a minimum temperature at or below 0°F, plus the count of days with at least 18 inches of snow on the ground; a day meeting both conditions counts twice. Conventional bands: below 50 mild, 50–80 moderate, 80–100 moderately severe, above 100 very severe. Read from the merged county-year panel, not recomputed here. CAVEAT: snow depth is an ERA5-Land grid-box average averaged again over a county, and that spatial averaging removes the local maxima an 18-inch cutoff is meant to catch, so across most of the eastern and midwestern deer range the index is driven almost entirely by its cold-day component.\par}
\clearpage
\begin{figure}[H]
\centering
\includegraphics[width=\textwidth]{\figTwo/prism_winter_severity_index_sensitivity.pdf}
\caption{The index under Kohn's 18-inch snow threshold and under a 12-inch alternative, on a shared scale.}
\end{figure}
\vspace{-6pt}
{\footnotesize\textit{Notes.} The same index with the snow-hazard component's threshold lowered from Kohn's 18 inches to 12 inches, on a shared colour scale. This is a robustness check, not an alternative definition of the Kohn index: it shows how much of the mapped variation is carried by the cold-day component alone once county-level spatial averaging has flattened the snow signal.\par}
\clearpage
\section{Anomalously warm winters}
\label{sec:anom}
Two thresholds are shown rather than one, and each map is paired with its own
time series. The construction matches \path{warm_winter_1sd} and
\path{warm_winter_2sd} in the merged county-year panel.
\par
\begin{figure}[H]
\centering
\includegraphics[width=\textwidth]{\figTwo/prism_anomalous_warm_winter_frequency_1sd.pdf}
\caption{Share of a county's winters more than one standard deviation above its own mean.}
\end{figure}
\vspace{-6pt}
{\footnotesize\textit{Notes.} Source: PRISM; complete DJF winters, 1982–2025. A winter is flagged when the county's mean winter temperature exceeds that county's own full-sample mean by more than 1 standard deviations — the identical construction to warm\_winter\_1sd in the merged county-year panel, so this map describes the regressor used in the collisions model. Colour scales differ between the 1 SD and 2 SD maps; 2 SD winters are far rarer.\par}
\clearpage
\begin{figure}[H]
\centering
\includegraphics[width=\textwidth]{\figTwo/prism_anomalous_warm_winter_time_series_1sd.pdf}
\caption{Share of counties flagged at one standard deviation, by winter year.}
\end{figure}
\vspace{-6pt}
{\footnotesize\textit{Notes.} Source: PRISM; complete DJF winters. Each bar is the share of counties whose winter mean exceeded that county's own full-sample mean by more than 1 standard deviations. Warm winters arrive as national episodes rather than as a steady drift, and they are spread across the sample rather than concentrated in recent years. Companion to the 1 SD frequency map.\par}
\clearpage
\begin{figure}[H]
\centering
\includegraphics[width=\textwidth]{\figTwo/prism_anomalous_warm_winter_frequency_2sd.pdf}
\caption{The same map at a two-standard-deviation threshold. Note the independent colour scale.}
\end{figure}
\vspace{-6pt}
{\footnotesize\textit{Notes.} Source: PRISM; complete DJF winters, 1982–2025. A winter is flagged when the county's mean winter temperature exceeds that county's own full-sample mean by more than 2 standard deviations — the identical construction to warm\_winter\_2sd in the merged county-year panel, so this map describes the regressor used in the collisions model. Colour scales differ between the 1 SD and 2 SD maps; 2 SD winters are far rarer.\par}
\clearpage
\begin{figure}[H]
\centering
\includegraphics[width=\textwidth]{\figTwo/prism_anomalous_warm_winter_time_series_2sd.pdf}
\caption{Share of counties flagged at two standard deviations, by winter year.}
\end{figure}
\vspace{-6pt}
{\footnotesize\textit{Notes.} Source: PRISM; complete DJF winters. Each bar is the share of counties whose winter mean exceeded that county's own full-sample mean by more than 2 standard deviations. Warm winters arrive as national episodes rather than as a steady drift, and they are spread across the sample rather than concentrated in recent years. Companion to the 2 SD frequency map.\par}
\clearpage
\section{Distributional shift}
\begin{figure}[H]
\centering
\includegraphics[width=\textwidth]{\figTwo/prism_national_winter_temperature_deviations_by_decade.pdf}
\caption{County-winter temperatures by period, after subtracting each county's own full-sample mean.}
\end{figure}
\vspace{-6pt}
{\footnotesize\textit{Notes.} Source: PRISM; all complete county-winters, each expressed as a deviation from its own county's full-sample mean. Black bars mark medians; violin width shows density. Removing persistent geography isolates the temporal shift. Periods contain unequal numbers of winters — 1980s: n=24,864; 1990s: n=31,080; 2000s: n=31,080; 2010s: n=31,080; 2020–2025: n=18,648. Estimand: the average county (counties weighted equally), not an area-, population- or road-weighted national average.\par}
\clearpage
\begin{figure}[H]
\centering
\includegraphics[width=\textwidth]{\figTwo/prism_national_winter_temperature_deviations_by_decade_area_weighted.pdf}
\caption{The same distributions with counties weighted by land area.}
\end{figure}
\vspace{-6pt}
{\footnotesize\textit{Notes.} Source: PRISM; the same county-winter deviations as the preceding exhibit, with each county weighted by its land area rather than counted once. Densities are weighted Gaussian kernel estimates and medians are weighted medians. Estimand: the average square kilometre of the contiguous United States (counties weighted by land area). Eyal ruled out population weighting (8/28); counties in the west are much larger than counties in the east, so the unweighted and area-weighted series answer different questions.\par}
\clearpage
\section{Winter warming relative to the other seasons}
Whether knowing how much a county warmed in summer tells us how much it warmed
in winter. If the seasons had moved together, these would lie on the 45-degree
line.
\par
\begin{figure}[H]
\centering
\includegraphics[width=\textwidth]{\figTwo/prism_cross_season_anomaly_correlation.pdf}
\caption{Winter temperature anomaly against the same year's spring, summer and fall anomalies.}
\end{figure}
\vspace{-6pt}
{\footnotesize\textit{Notes.} Source: PRISM; one observation per county-year, standardised within county and season against that county-season's own 1982–2025 mean and standard deviation — the same construction as the warm-winter exhibits. Dots are the means of 20 equal-count bins; the orange line is the pooled OLS fit and the grey dashed line is the 45° line. Points on the 45° line would mean a season's anomaly is fully informative about the same year's winter anomaly.\par}
\clearpage
\begin{figure}[H]
\centering
\includegraphics[width=\textwidth]{\figTwo/prism_cross_season_trend_correlation.pdf}
\caption{County winter warming trend against that county's spring, summer and fall trends.}
\end{figure}
\vspace{-6pt}
{\footnotesize\textit{Notes.} Source: PRISM; one observation per county — the OLS trend in that county's seasonal mean temperature, 1982–2025, in °C per decade. Grey points are counties, coloured dots are the means of 20 equal-count bins, orange is the pooled fit and the grey dashed line is the 45° line. Counties below the 45° line have warmed LESS in winter than in the season on the horizontal axis.\par}
\clearpage
\section{Robustness: PRISM versus ERA5}
\begin{figure}[H]
\centering
\includegraphics[width=\textwidth]{\figTwo/prism_era5_national_winter_temperature_comparison.pdf}
\caption{National winter temperature from both products. PRISM is the analysis dataset.}
\end{figure}
\vspace{-6pt}
{\footnotesize\textit{Notes.} ERA5 minus PRISM: mean +0.39°C (range +0.07 to +0.73); correlation 0.993. The two products track each other closely and the gap does not widen or close systematically over the sample, which is the property that matters for robustness. PRISM is the analysis dataset; ERA5 is a robustness check only. Estimand: the average county (counties weighted equally), not an area-, population- or road-weighted national average.\par}
\clearpage
\begin{figure}[H]
\centering
\includegraphics[width=\textwidth]{\figTwo/prism_era5_national_winter_temperature_comparison_area_weighted.pdf}
\caption{The same comparison with counties weighted by land area.}
\end{figure}
\vspace{-6pt}
{\footnotesize\textit{Notes.} The same PRISM–ERA5 comparison with counties weighted by land area: mean offset +0.01°C, correlation 0.985. Estimand: the average square kilometre of the contiguous United States (counties weighted by land area). Eyal ruled out population weighting (8/28); counties in the west are much larger than counties in the east, so the unweighted and area-weighted series answer different questions.\par}
\clearpage
\section{Decisions requested}
\label{sec:decisions}
\subsection*{Mean winter temperature: two constructions}
\textbf{Current choice.} This notebook takes a DAY-COUNT-WEIGHTED mean of the three monthly means and keeps only winters complete on both month and day counts. build\_main\_data\_county\_year.do line 512 takes a simple average of the three monthly means.\par
\textbf{Issue.} The gap is hundredths of a degree, but it is enough to move a county-winter across a sigma threshold, so the anomaly counts here and in the merged panel will not match exactly. Harmonise the merge script to the day-weighted mean, or accept the difference and document it? Table 1 has already been run off the current merge script.\par
\textbf{Status.} DECIDE
\subsection*{Anomaly definition}
\textbf{Current choice.} Raw winter mean above the county's own full-sample mean by 1 or 2 SD — identical to warm\_winter\_1sd / warm\_winter\_2sd in SECTION 7 of the merge script.\par
\textbf{Issue.} Changed in this pass. The v1 exhibits used a DETRENDED residual above 1.5 SD with a 30-winter filter, which described a different object from the regressor. The detrended version is retained as a Tier 1 diagnostic. Flagging in case the detrended construction should be the headline instead.\par
\textbf{Status.} CHANGED — confirm
\subsection*{Geographic weighting}
\textbf{Current choice.} Unweighted is primary; an area-weighted companion ships for each exhibit that pools counties into a national number.\par
\textbf{Issue.} Population weighting was ruled out on 8/28. For the collisions stage the relevant exposure weight may be vehicle-miles travelled or road miles rather than land area.\par
\textbf{Status.} CONFIRM
\subsection*{Comparison periods}
\textbf{Current choice.} 1982–1991 versus 2016–2025\par
\textbf{Issue.} Symmetric ten-winter period means, not 1982-versus-2025 endpoints.\par
\textbf{Status.} CONFIRM
\subsection*{Region scheme for the pooled state figure}
\textbf{Current choice.} EPA climate-impact regions, assigned whole states (both variants ship).\par
\textbf{Issue.} EPA Level I ecoregions cut across state lines and would need a county-level spatial join. Pick the plain or the region-coloured version for the paper.\par
\textbf{Status.} PICK ONE
\subsection*{Winter Severity Index exhibit}
\textbf{Current choice.} Read from the merged county-year panel; Kohn 1975 / WI DNR definition.\par
\textbf{Issue.} County-level averaging of ERA5-Land snow depth under-detects the 18-inch threshold, so across most of the eastern deer range the index is effectively its cold-day component. The 12-inch sensitivity panel shows this.\par
\textbf{Status.} DOCUMENTED
\subsection*{Extreme-cold change colour scale}
\textbf{Current choice.} Reversed relative to the temperature-change map: red marks FEWER cold days.\par
\textbf{Issue.} Deliberate — for a count of cold days, a decrease is the warming signal.\par
\textbf{Status.} DOCUMENTED
\clearpage
\appendix
\section{State-level winter-temperature trends}
\label{app:states}
Per-state annual means with state-specific linear trends and the standard error
of each slope. Reproduced in full; page count follows the run. Y-axis ranges
differ by panel deliberately -- a common range across states whose mean winter
temperature spans roughly 40\textdegree C would flatten every trend to invisibility.
\includepdf[pages=-]{\figTwo/prism_state_long_run_warming.pdf}
\section{Tier 1 QA artefacts (not reproduced)}
The following internal QA exhibits and tables were produced by the same run and
are available on request. This includes the pooled-level distribution exhibit,
which was demoted from the report set, and the detrended warm-anomaly diagnostic.
\begin{itemize}[leftmargin=*,itemsep=1pt]
\item \path{figures/weather/tier1/prism_interannual_variability_choropleth.pdf}
\item \path{figures/weather/tier1/prism_national_trend_days_below_freezing_32f.pdf}
\item \path{figures/weather/tier1/prism_national_trend_days_extremely_cold.pdf}
\item \path{figures/weather/tier1/prism_national_trend_freeze_thaw_days.pdf}
\item \path{figures/weather/tier1/prism_county_spaghetti_mean_temp_c.pdf}
\item \path{figures/weather/tier1/prism_state_trends_freeze_thaw_days.pdf}
\item \path{figures/weather/tier1/prism_state_decade_distributions_mean_temp_c.pdf}
\item \path{figures/weather/tier1/prism_winter_temperature_distributions_by_decade_tier1.pdf}
\item \path{figures/weather/tier1/prism_state_trends_mean_temp_c.pdf}
\item \path{figures/weather/tier1/prism_state_trends_days_below_freezing_32f.pdf}
\item \path{figures/weather/tier1/prism_national_trend_mean_temp_c.pdf}
\item \path{figures/weather/tier1/prism_state_trends_days_extremely_cold.pdf}
\item \path{tables/weather/tier1/prism_tier1_qa_findings.csv}
\item \path{tables/weather/tier1/prism_county_winter_warm_anomalies_detrended_diagnostic.csv}
\item \path{tables/weather/tier1/prism_summary_by_county_review_flags.csv}
\item \path{tables/weather/tier1/prism_summary_pooled.csv}
\item \path{tables/weather/tier1/prism_interannual_variability_by_state.csv}
\item \path{tables/weather/tier1/prism_summary_by_calendar_month.csv}
\item \path{tables/weather/tier1/prism_tier1_coverage_integrity.csv}
\item \path{tables/weather/tier1/prism_summary_by_county.csv}
\item \path{tables/weather/tier1/prism_county_interannual_variability.csv}
\end{itemize}
\end{document}
